\section{Loan Lifecycle}
\label{sec:lifecycle}

\subsection{Origination}

Origination follows the order book flow (Section~\ref{sec:architecture}): a borrower posts a request, solvers respond with quotes, and the accepted quote is committed into an epoch's settlement tree and executed atomically. Five origination pipelines exist, each assembled as an atomic hook chain:

\textbf{Vault-funded.} The vault transfers USDC principal directly to the borrower (or to the margin agent for Prime Trade positions). The collateral manager holds collateral without supplying it to an external venue; no hedging buffers are reserved.

\textbf{Externally-sourced.} The hook chain deposits collateral into the collateral manager, supplies it to the external venue, borrows USDC against it, and transfers the USDC to the borrower. A rate buffer is escrowed to hedge the mismatch between the borrower's fixed rate and the venue's variable rate.

\textbf{Externally-sourced with collateral swap.} For long-tail collateral not listed on the external venue (e.g.\ LDO, UNI): a swap hook converts the collateral to a venue-listed intermediate (e.g.\ LDO $\to$ WETH), then the chain proceeds as above. Both a rate buffer and an option buffer are reserved; the option buffer covers the round-trip swap cost at close.

\textbf{Hyperliquid margin.} For Prime Trade positions, the hook chain follows the externally-sourced pipeline but instead of transferring USDC to the borrower, a bridge hook burns the USDC via CCTP~\cite{cctp} to a per-loan margin agent on HyperEVM, where a delivery hook deposits it as margin (Section~\ref{sec:crosschain}). Collateral never leaves Ethereum.

\textbf{Hyperliquid margin with collateral swap.} Combines the collateral swap with Hyperliquid delivery: swap hook, option buffer, external venue borrow, then CCTP bridge to the margin agent.

\begin{figure}[ht]
  \centering
  \begin{tikzpicture}[
      box/.style={draw, thick, rounded corners=3pt, fill=gray!8,
      minimum height=1.2cm, minimum width=2.8cm, align=center, font=\small\sffamily},
      arr/.style={-{Stealth[length=8pt]}, thick},
      lab/.style={font=\small\sffamily, fill=white, inner sep=3pt}
    ]
    \node[box] (B)   at (0,0)      {Borrower};
    \node[box] (CDP) at (0,-2.5)   {CDP};
    \node[box] (AV)  at (0,-5)     {Ext.\ Venue};
    \node[box] (V)   at (4.5,-2.5) {Vault};

    \draw[arr] ([xshift=-10pt]B.south) -- node[lab,left]{1.\ WETH} ([xshift=-10pt]CDP.north);
    \draw[arr] ([xshift=10pt]CDP.north) -- node[lab,right]{4.\ USDC} ([xshift=10pt]B.south);
    \draw[arr] ([xshift=-10pt]CDP.south) -- node[lab,left]{2.\ supply} ([xshift=-10pt]AV.north);
    \draw[arr] ([xshift=10pt]AV.north) -- node[lab,right]{3.\ USDC} ([xshift=10pt]CDP.south);
    \draw[arr] (CDP) -- node[lab,above]{5.\ debt tokens} (V);
  \end{tikzpicture}
  \caption{Origination, externally-sourced pipeline. Steps 1--4 execute as a single atomic hook chain. Debt tokens (step~5) are minted to the vault; the rate buffer is escrowed. Vault-funded origination omits the external venue and transfers principal directly.}
  \label{fig:origination}
\end{figure}

\subsection{Refinancing}

A borrower can refinance through the order book rather than repaying and re-borrowing manually: the borrower posts a refinancing intent, solvers respond with quotes, and the matching engine treats these identically to origination quotes. On match, the existing loan closes and a new loan originates atomically within the same epoch. When the Harvest Vault is creditor on both loans, the refinancing is a pure accounting event.

\subsection{Close and Liquidation}

Repayment follows a two-phase process. The borrower deposits USDC into the issued debt contract, increasing its redeemable reserve; after the next interest accrual is verified (Section~\ref{sec:tokens}), the vault redeems its debt tokens and routes the USDC to the CDP for variable-rate debt repayment on the external venue.

On voluntary close, the platform unwinds the external position atomically via a close hook chain: repays the variable debt, withdraws the collateral, and returns it to the borrower (with a reverse swap for collateral-swap loans). Any hedging surplus stays in the vault as LP yield. Vault-funded loans skip the external unwind: the borrower repays, and collateral is released directly.

Liquidation spans two epochs. In the first, a proof attesting the liquidatable state is committed as a leaf in the epoch's liquidation root; detection is identity-agnostic, requiring only the mathematical fact that barriers have been breached. In the second, the liquidator identifies itself on-chain, repays the outstanding debt via the debt token contract, and seizes collateral from the CDP at a risk-engine-determined price. Separating detection from seizure ensures collateral moves only after debt repayment is recorded on-chain.
